\documentclass[aspectratio=169]{beamer}
\usetheme{metropolis}
\metroset{sectionpage=none, numbering=counter, progressbar=frametitle}
\usepackage{amsmath, amssymb}
\usepackage[T1]{fontenc}

\title{Domain-Partition Algorithm for Turbomachinery Blades}
\subtitle{Slide 2: Unwrapped Hub Surface}
\date{\today}
\author{}
\institute{}

\begin{document}

\begin{frame}{Unwrapped Hub Surface --- Cylindrical Mapping}

\begin{center}
\Large
\[
\theta = \arctan\!\left(\frac{y}{x}\right), \quad
s = r \cdot \theta, \quad
t = z
\]
\end{center}

\vspace{0.5em}
\begin{itemize}
    \item \textbf{Isometric mapping:} distortion-free, preserves triangulation
    \item \textbf{Domain becomes rectangular} with a blade-shaped hole
    \item \textbf{Periodic $\theta$-boundaries} (pitch = $\pi/4$ for 8 blades)
\end{itemize}

\vspace{0.8em}
\small
\textit{Code:} \texttt{scripts/stage1/unwrap\_surface.py:8--10}

\note{
Speaker notes:
The mapping theta = atan2(y, x), s = r*theta, t = z is isometric because the cylinder has constant radius. This means the 2D triangulation is identical to the 3D one, only vertex coordinates change. The resulting domain is rectangular with a blade-shaped hole, exactly the 2D domain-partition case. For T1_9 with 8 blades, the periodic pitch is pi/4.
}

\end{frame}

\end{document}
